\section{幺模群}

设 n 是一个正整数。我们用 \(\mathbf{SL}_{n}(\mathrm{D})\) 表示群同态 \(\det: \mathbf{GL}_{n}(\mathrm{D}) \to \mathrm{D}_{\mathrm{ab}}^{*}\) 的核，并称它为幺模群或特殊线性群（当体 D 交换时，参见 III, §8, No.~9, 第 537 页）。

定理 1. --- 设 \(n \geqslant 2\)。

\begin{enumerate}
\def\labelenumi{\alph{enumi})}
\item
  子群 \(\mathbf{SL}_n(\mathrm{D})\)（\(\mathbf{GL}_n(\mathrm{D})\) 的）由诸矩阵 \(B_{ij}(\lambda)\) 生成，其中 \(i\) 与 \(j\) 是区间 \([1,n]\) 中两个不同的整数，而 \(\lambda\) 跑遍 D。
\item
  设 \(n \geqslant 3\) 或 \(\operatorname{Card}(\mathrm{D}) \geqslant 3\)。\(\mathbf{GL}_n(\mathrm{D})\) 的导群等于 \(\mathbf{SL}_n(\mathrm{D})\)。
\item
  设 \(n \geqslant 3\) 或 \(\operatorname{Card}(\mathrm{D}) \geqslant 4\)。\(\mathbf{SL}_n(\mathrm{D})\) 的导群等于 \(\mathbf{SL}_n(\mathrm{D})\)。
\end{enumerate}

\begin{enumerate}
\def\labelenumi{\Alph{enumi})}
\tightlist
\item
  用 T 表示 \(\mathbf{GL}_{n}(\mathrm{D})\) 的由诸矩阵 \(B_{ij}(\lambda)\) 生成的子群。由 VIII, 第 454 页的例 1，我们有 \(\det(B_{ij}(\lambda)) = 1\)，从而 \(\mathrm{T} \subset \mathbf{SL}_{n}(\mathrm{D})\)。为证明这两个群相等，于是由同前所引的例 3，只需证明每个形如 \(\operatorname{diag}(1, \ldots, 1, d)\) 且 \(\pi(d) = 1\) 的矩阵属于 T。矩阵 \(\operatorname{diag}(1, \ldots, 1, d)\) 属于同态 \(U \mapsto \begin{pmatrix} I_{n-2} & 0 \\ 0 & U \end{pmatrix}\)（从 \(\mathbf{GL}_{2}(\mathrm{D})\) 到 \(\mathbf{GL}_{n}(\mathrm{D})\)）的像；因此只需考虑 n = 2 的情形。由于 \(\pi\) 的核是 \(D^{*}\) 的导群，我们可以假设 \(d = uvu^{-1}v^{-1}\)，其中 u、v 属于 \(D^{*}\)。于是我们的断言由等式
\end{enumerate}

\[
\left( \begin{array}{c c} 1 & 0 \\ 0 & d \end{array} \right) = \left( \begin{array}{c c} u ^ {- 1} & 0 \\ 0 & u \end{array} \right) \left( \begin{array}{c c} v ^ {- 1} & 0 \\ 0 & v \end{array} \right) \left( \begin{array}{c c} v u & 0 \\ 0 & u ^ {- 1} v ^ {- 1} \end{array} \right)\tag{28}
\]

与

\[
\left( \begin{array}{c c} s & 0 \\ 0 & s ^ {- 1} \end{array} \right) = \left( \begin{array}{c c} 1 & s \\ 0 & 1 \end{array} \right) \left( \begin{array}{c c} 1 & 0 \\ 1 - s ^ {- 1} & 1 \end{array} \right) \left( \begin{array}{c c} 1 & - 1 \\ 0 & 1 \end{array} \right) \left( \begin{array}{c c} 1 & 0 \\ 1 - s & 1 \end{array} \right)\tag{29}
\]

（对 \(s \in D^{*}\)）得出。

\begin{enumerate}
\def\labelenumi{\Alph{enumi})}
\setcounter{enumi}{1}
\tightlist
\item
  按构造，\(\mathbf{SL}_{n}(\mathrm{D})\) 是一个从 \(\mathbf{GL}_{n}(\mathrm{D})\) 到交换群的同态的核，所以它包含导群 \((\mathbf{GL}_{n}(\mathrm{D}),\mathbf{GL}_{n}(\mathrm{D}))\)（\(\mathbf{GL}_{n}(\mathrm{D})\) 的）。由上述，我们有
\end{enumerate}

\[
\mathbf {S L} _ {n} (\mathrm{D}) \supset (\mathbf {G L} _ {n} (\mathrm{D}), \mathbf {G L} _ {n} (\mathrm{D})) \supset (\mathbf {S L} _ {n} (\mathrm{D}), \mathbf {S L} _ {n} (\mathrm{D})).
\]

为证明 c)，只需证明诸矩阵 \(B_{ij}(\lambda)\) 是 \(\mathbf{SL}_n(\mathrm{D})\) 中的换位子。

设 \(n \geqslant 3\)。若 \(i, j, k\) 是区间 \([1, n]\) 中三个不同的整数，且 \(\mu, \nu\) 是 \(D\) 的元素，则我们有

\[
B _ {i j} (\mu \nu) = B _ {i k} (\mu) ^ {- 1} B _ {k j} (\nu) ^ {- 1} B _ {i k} (\mu) B _ {k j} (\nu).\tag{30}
\]

取 \(\mu = 1\) 与 \(\nu = \lambda\)，我们看出矩阵 \(B_{ij}(\lambda)\) 是 \(\mathbf{SL}_n(\mathrm{D})\) 的元素的换位子。

现在设 \(n = 2\)。设 \(u\) 与 \(v\) 是 \(D\) 的元素且 \(u \neq 0\)。我们有关系式

\[
\left( \begin{array}{c c} 1 & v - u v u \\ 0 & 1 \end{array} \right) = \left( \begin{array}{c c} u & 0 \\ 0 & u ^ {- 1} \end{array} \right) \left( \begin{array}{c c} 1 & - v \\ 0 & 1 \end{array} \right) \left( \begin{array}{c c} u ^ {- 1} & 0 \\ 0 & u \end{array} \right) \left( \begin{array}{c c} 1 & v \\ 0 & 1 \end{array} \right)\tag{31}
\]

与

\[
\left( \begin{array}{c c} 1 & 0 \\ v - u v u & 1 \end{array} \right) = \left( \begin{array}{c c} u ^ {- 1} & 0 \\ 0 & u \end{array} \right) \left( \begin{array}{c c} 1 & 0 \\ - v & 1 \end{array} \right) \left( \begin{array}{c c} u & 0 \\ 0 & u ^ {- 1} \end{array} \right) \left( \begin{array}{c c} 1 & 0 \\ v & 1 \end{array} \right).\tag{32}
\]

我们有 \(\det\left(\begin{matrix}u & 0 \\ 0 & u^{-1}\end{matrix}\right)=1\)，所以矩阵 \(B_{12}(v-uvu)\) 与 \(B_{21}(v-uvu)\) 是 \(\mathbf{SL}_{n}(\mathrm{D})\) 的元素的换位子。

设体 D 至少有四个元素。设 \(\lambda\) 是 D 的一个元素。若 \(\lambda\) 等于 0、1 或 -1，则在 \(D-\{0,1,-1\}\) 中任意选取一个元素 u；否则，置 \(u=\lambda\)。在这两种情形下，u 都是 D 的非零元素，它与 \(\lambda\) 交换，且我们有 \(u^{2}\neq1\)。置 \(v=\lambda(1-u^{2})^{-1}\)。我们有 uv=vu，从而 \(v-uvu=v(1-u^{2})=\lambda\)。于是由关系式 (31) 与 (32) 得出，矩阵 \(B_{12}(\lambda)\) 与 \(B_{21}(\lambda)\) 是 \(\mathbf{SL}_{n}(\mathrm{D})\) 中的换位子。这就完成了 c) 的证明。

\begin{enumerate}
\def\labelenumi{\Alph{enumi})}
\setcounter{enumi}{2}
\tightlist
\item
  只剩下证明当 D 有三个元素时 \(\mathbf{SL}_{2}(\mathrm{D})\) 是 \(\mathbf{GL}_{2}(\mathrm{D})\) 的导群。由 A)，群 \(\mathbf{SL}_{2}(\mathrm{D})\) 由矩阵 \(B_{12}(1)=(\begin{matrix}{1}&{1}\\ {0}&{1}\end{matrix})\) 与 \(B_{21}(1)=(\begin{matrix}{1}&{0}\\ {1}&{1}\end{matrix})\) 生成，且这些矩阵是 \(\mathbf{GL}_{2}(\mathrm{D})\) 的元素的换位子，因为我们有 \(B_{21}(1)={}^{t}B_{12}(1)\)，且
\end{enumerate}

\[
\left( \begin{array}{c c} 1 & 1 \\ 0 & 1 \end{array} \right) = \left( \begin{array}{c c} - 1 & 0 \\ 0 & 1 \end{array} \right) \left( \begin{array}{c c} 1 & 1 \\ 0 & 1 \end{array} \right) \left( \begin{array}{c c} - 1 & 0 \\ 0 & 1 \end{array} \right) ^ {- 1} \left( \begin{array}{c c} 1 & 1 \\ 0 & 1 \end{array} \right) ^ {- 1}.\tag{33}
\]

注. --- 若 D 是一个只有两个元素的体，则 \(\mathbf{GL}_{2}(\mathrm{D})\) 等于 \(\mathbf{SL}_{2}(\mathrm{D})\)，它是一个 6 阶群，其导群的阶为 3。若 D 是一个只有三个元素的体，则群 \(\mathbf{SL}_{2}(\mathrm{D})\) 的阶为 24，其导群的阶为 8。设 \(n \geqslant 2\)；除前两种情形外，\(\mathbf{SL}_{n}(\mathrm{D})\) 的每个异于 \(\mathbf{SL}_{n}(\mathrm{D})\) 的正规子群都包含在中心 \(\mathrm{Z}(\mathbf{SL}_{n}(\mathrm{D}))\)（\(\mathbf{SL}_{n}(\mathrm{D})\) 的）中（II, §10, 第 421 页，习题 13 与 14）。这时群 \(\mathbf{SL}_{n}(\mathrm{D})/\mathrm{Z}(\mathbf{SL}_{n}(\mathrm{D}))\) 是单群。

设 V 是体 D 上的一个右向量空间，维数 n 有限。我们用 \(\mathrm{SL}(V)\) 表示同态 \(\det: \mathbf{GL}(V) \to D_{ab}^{*}\) 的核，并称之为 V 的幺模群。选取 V 的一个基就能把这个群与 \(\mathbf{SL}_{n}(D)\) 等同起来。

我们称 V 的一个自同构 u 为平延，如果存在 V 的一个向量 v 与 V 上的一个线性型 \(\varphi\)，使得 \(\varphi(v)=0\) 且对 V 中的每个 x 有 \(u(x)=x+v\varphi(x)\)。当 \(n\geqslant2\) 时，u 是平延当且仅当存在 V 的一个基，u 关于这个基的矩阵形如 \(B_{ij}(\lambda)\)。定理 1 蕴含下面的推论。

推论. --- 设 V 是体 D 上的一个右向量空间，维数 \(n \geqslant 2\)。

\begin{enumerate}
\def\labelenumi{\alph{enumi})}
\item
  子群 \(\mathbf{SL}(\mathrm{V})\)（\(\mathbf{GL}(\mathrm{V})\) 的）由诸平延生成。
\item
  子群 \(\mathbf{SL}(\mathrm{V})\) 是 \(\mathbf{GL}(\mathrm{V})\) 的导群，除非 \(n = 2\) 且 D 只有两个元素。
\item
  除非 n = 2 且 D 只有两个或三个元素，群 \(\mathbf{SL}(V)\) 等于它的导群。
\end{enumerate}

